% ============================================================================
% SUPPLEMENTARY MATERIAL: Mathematical Formulations
% Saved from ml_chemistry_template.tex (originally lines 709-931)
% These formulations provide detailed mathematical descriptions of the
% RNN architectures used in the paper.
% ============================================================================

\documentclass[11pt,a4paper]{article}
\usepackage[utf8]{inputenc}
\usepackage[T1]{fontenc}
\usepackage{amsmath,amssymb,amsfonts}
\usepackage{graphicx}
\usepackage{hyperref}
\usepackage[margin=1in]{geometry}
\usepackage{tikz}
\usetikzlibrary{arrows.meta, positioning, shapes, calc}

\newcommand{\vect}[1]{\mathbf{#1}}
\newcommand{\matr}[1]{\mathbf{#1}}

\title{Supplementary Material: Mathematical Formulations for Sequential RNN Deep Networks}
\author{}
\date{}

\begin{document}
\maketitle

\section{Sequence-based deep learning models for UV absorption maximum prediction}
\label{sec:seq_uv_models}

In this work, we develop sequence-based deep learning models for predicting the UV--visible absorption maximum, $\lambda_{\max}$, of small organic UV absorbers from their SMILES representations. The UV absorption maximum is a key photophysical descriptor for applications such as sunscreens, UV-protective coatings, and UV-stabilizing additives, and is closely related to the underlying electronic structure and conjugation pattern of the chromophore. Conventional quantitative structure--property relationship (QSPR) models for $\lambda_{\max}$ rely on handcrafted descriptors (e.g., conjugation length, Hammett parameters, quantum-chemical features), which can be expensive to compute and may generalize poorly to new chemical scaffolds.

By contrast, sequence-based neural models treat SMILES strings as one-dimensional molecular ``languages'' and learn task-specific representations directly from the raw sequence. In particular, we employ gated recurrent neural networks (GRUs/LSTMs) over SMILES sequences to encode molecular structure into a fixed-size vector, which is then used to regress $\lambda_{\max}$ in nanometers. This section describes the input representation, recurrent architectures, and model design choices adopted for UV absorption maximum prediction.

\subsection{Sequence representation and notation}
\label{subsec:notation}

Let a molecule be represented by a SMILES string $s$, which we treat as a sequence of discrete tokens:
\begin{equation}
  s = (c_1, c_2, \dots, c_T), \quad c_t \in \mathcal{V},
\end{equation}
where $T \in \mathbb{N}$ is the sequence length, and $\mathcal{V}$ is a finite vocabulary of SMILES tokens. Each token $c_t$ is mapped to a numerical input vector $x_t \in \mathbb{R}^{d_{\mathrm{emb}}}$ via an embedding layer (Section~\ref{subsec:input_representation}), producing an embedded sequence
\begin{equation}
  (x_1, x_2, \dots, x_T), \quad x_t \in \mathbb{R}^{d_{\mathrm{emb}}}.
\end{equation}

A recurrent encoder then computes hidden states
\begin{equation}
  h_t = f_{\mathrm{RNN}}(x_t, h_{t-1}; \theta), \quad h_t \in \mathbb{R}^{d_h},
\end{equation}
where $d_h$ is the hidden dimension and $\theta$ denotes the trainable parameters of the recurrent cell. Depending on the architecture, we may also maintain a cell state $c_t \in \mathbb{R}^{d_h}$ (for LSTMs). For bidirectional encoders, we obtain forward and backward hidden states at each position $t$, denoted by $\overrightarrow{h}_t$ and $\overleftarrow{h}_t$, which are concatenated to form a context-aware token representation
\begin{equation}
  h_t = [\,\overrightarrow{h}_t; \overleftarrow{h}_t\,] \in \mathbb{R}^{2d_h}.
\end{equation}

The sequence of hidden states $(h_1, \dots, h_T)$ is subsequently aggregated into a fixed-size molecular representation $h_{\mathrm{mol}}$ by a pooling layer, which is then passed to dense layers to predict $\lambda_{\max}$.

\subsection{Input representation: SMILES tokenization, one-hot encoding, and embeddings}
\label{subsec:input_representation}

\paragraph{Tokenization.}
We adopt \emph{character-level} tokenization for SMILES strings due to its simplicity and robustness. The vocabulary $\mathcal{V}$ contains all SMILES characters observed in the dataset (e.g., atom symbols such as \texttt{C}, \texttt{N}, \texttt{O}; bond symbols such as \texttt{=}, \texttt{\#}; branching symbols \texttt{(}, \texttt{)}; ring indices; and brackets), augmented with special tokens:
\begin{itemize}
  \item \verb|<bos>|: beginning-of-sequence marker,
  \item \verb|<eos>|: end-of-sequence marker,
  \item \verb|<pad>|: padding token for batched training,
  \item \verb|<unk>|: unknown token (optional, for rare/unseen characters).
\end{itemize}
Character-level tokenization avoids the need for dataset-specific parsing rules and works seamlessly with randomized SMILES augmentation. It also ensures that the same tokenization scheme can be applied consistently across different UV absorber datasets.

\paragraph{One-hot encoding.}
Each token $c_t \in \mathcal{V}$ is mapped to an integer index
\begin{equation}
  k_t = \text{index}(c_t) \in \{1, 2, \dots, |\mathcal{V}|\},
\end{equation}
and then to a one-hot vector $o_t \in \{0,1\}^{|\mathcal{V}|}$ defined by
\begin{equation}
  (o_t)_j =
  \begin{cases}
    1, & \text{if } j = k_t, \\
    0, & \text{otherwise}.
  \end{cases}
\end{equation}
One-hot vectors are high-dimensional, sparse, and do not encode any inherent similarity between tokens (all distinct tokens are equidistant in Euclidean space). Consequently, they serve only as an intermediate representation and are not fed directly into the recurrent layers.

\paragraph{Embedding layer.}
To obtain compact and trainable token features, we introduce a learnable embedding matrix
\begin{equation}
  E \in \mathbb{R}^{|\mathcal{V}| \times d_{\mathrm{emb}}},
\end{equation}
with embedding dimension fixed to $d_{\mathrm{emb}} = 50$ in all experiments. The embedded input at time step $t$ is given by
\begin{equation}
  x_t = E^{\top} o_t \in \mathbb{R}^{d_{\mathrm{emb}}} = \mathbb{R}^{50},
\end{equation}
which is equivalent in practice to selecting the row of $E$ corresponding to token $c_t$. The embedding layer thus implements a mapping
\begin{equation}
  \text{Embed} : \mathcal{V} \to \mathbb{R}^{50}, \quad c_t \mapsto x_t.
\end{equation}
The parameters of $E$ are initialized randomly and learned jointly with the recurrent encoder and prediction head via backpropagation. Over the course of training, the model discovers a task-specific geometry in embedding space: tokens that play similar roles in SMILES (e.g., common heavy atoms or paired structural delimiters) may acquire similar embeddings, while tokens with distinct functional roles are separated. After embedding, each SMILES string is represented as
\begin{equation}
  (x_1, x_2, \dots, x_T), \quad x_t \in \mathbb{R}^{50},
\end{equation}
which serves as the input to the recurrent encoder.

\subsection{Recurrent architectures: RNN, LSTM, and GRU}
\label{subsec:rnn_architectures}

\paragraph{Vanilla RNN.}
The Elman RNN defines the hidden state update as
\begin{equation}
  h_t = \phi\left( W_{xh} x_t + W_{hh} h_{t-1} + b_h \right),
\end{equation}
where $W_{xh} \in \mathbb{R}^{d_h \times d_{\mathrm{emb}}}$, $W_{hh} \in \mathbb{R}^{d_h \times d_h}$, $b_h \in \mathbb{R}^{d_h}$, and $\phi(\cdot)$ is a nonlinear activation (typically $\tanh$). Although conceptually simple, vanilla RNNs suffer from vanishing and exploding gradients for longer sequences, making it difficult to capture long-range dependencies in SMILES strings (such as interactions between distant segments of a conjugated chromophore). As a result, vanilla RNNs are considered only as a conceptual baseline and are not used as the primary architecture.

\paragraph{Long Short-Term Memory (LSTM).}
LSTMs augment the RNN with a cell state $c_t$ and gating mechanisms to better preserve information over long time spans. At each time step, input, forget, and output gates $i_t, f_t, o_t \in (0,1)^{d_h}$, together with a candidate cell state $\tilde{c}_t$, are computed as
\begin{align}
  i_t &= \sigma(W_{xi} x_t + W_{hi} h_{t-1} + b_i), \\
  f_t &= \sigma(W_{xf} x_t + W_{hf} h_{t-1} + b_f), \\
  o_t &= \sigma(W_{xo} x_t + W_{ho} h_{t-1} + b_o), \\
  \tilde{c}_t &= \tanh(W_{xc} x_t + W_{hc} h_{t-1} + b_c),
\end{align}
followed by the cell and hidden state updates
\begin{align}
  c_t &= f_t \odot c_{t-1} + i_t \odot \tilde{c}_t, \\
  h_t &= o_t \odot \tanh(c_t),
\end{align}
where $\sigma(\cdot)$ denotes the logistic sigmoid function and $\odot$ is the element-wise product. The cell state $c_t$ provides a controlled pathway for information to flow across time steps, which is important for modeling non-local structural relationships in SMILES (e.g., conjugated segments or ring systems spanning many tokens). In this work, LSTMs are used as alternative recurrent cells in ablation studies.

\paragraph{Gated Recurrent Unit (GRU).}
GRUs provide a simpler gated architecture with a single hidden state and no separate cell state. At each time step, update and reset gates $z_t, r_t \in (0,1)^{d_h}$ are computed as
\begin{align}
  z_t &= \sigma(W_{xz} x_t + W_{hz} h_{t-1} + b_z), \\
  r_t &= \sigma(W_{xr} x_t + W_{hr} h_{t-1} + b_r),
\end{align}
and the candidate hidden state $\tilde{h}_t$ as
\begin{equation}
  \tilde{h}_t = \tanh \left( W_{xh} x_t + W_{hh} (r_t \odot h_{t-1}) + b_h \right).
\end{equation}
The new hidden state is then
\begin{equation}
  h_t = (1 - z_t) \odot h_{t-1} + z_t \odot \tilde{h}_t.
\end{equation}
GRUs retain the benefits of gating while using fewer parameters than LSTMs. Empirically, GRUs and LSTMs often achieve comparable performance when appropriately tuned, while GRUs can be slightly more efficient to train, which is advantageous in the moderate data regimes typical of curated UV absorber datasets.

\paragraph{Bidirectional and stacked encoders.}
To incorporate both left and right SMILES context, we employ \emph{bidirectional} GRU encoders. A single bidirectional GRU (BiGRU) layer computes forward and backward hidden states,
\begin{equation}
  \overrightarrow{h}_t = f_{\mathrm{GRU}}(x_t, \overrightarrow{h}_{t-1}), \quad
  \overleftarrow{h}_t = f_{\mathrm{GRU}}(x_t, \overleftarrow{h}_{t+1}),
\end{equation}
which are concatenated,
\begin{equation}
  h_t^{(1)} = [\,\overrightarrow{h}_t; \overleftarrow{h}_t\,] \in \mathbb{R}^{2d_h}.
\end{equation}
In our architecture, we set $d_h = 128$ for each unidirectional GRU, so $h_t^{(1)} \in \mathbb{R}^{256}$. We stack \textbf{two} such bidirectional GRU layers. The second bidirectional layer takes $h_t^{(1)}$ as input and produces
\begin{equation}
  h_t^{(2)} = [\,\overrightarrow{h}_t^{(2)}; \overleftarrow{h}_t^{(2)}\,] \in \mathbb{R}^{256},
\end{equation}
which forms the final sequence of token-level representations used for pooling.

\subsection{Model design choices for $\lambda_{\max}$ prediction}
\label{subsec:design_choices}

\paragraph{Recurrent encoder configuration.}
The main model consists of:
\begin{itemize}
  \item an embedding layer with dimension $d_{\mathrm{emb}} = 50$,
  \item two stacked bidirectional GRU layers, each with $d_h = 128$ units per direction (yielding 256-dimensional token representations after each layer),
  \item a global average pooling layer,
  \item a dense hidden layer with 128 units and ReLU activation,
  \item a dropout layer with dropout rate $0.2$,
  \item and a final dense layer with linear activation producing a scalar $\hat{y}$ for $\lambda_{\max}$.
\end{itemize}
GRUs are chosen as the primary recurrent cell due to their favorable trade-off between expressivity and parameter efficiency. LSTM-based variants with matching hidden size and depth are implemented for ablation experiments, confirming that the GRU-based encoder provides comparable or better performance with slightly lower computational cost.

\paragraph{Sequence-to-vector pooling (attention).}
Given the sequence of hidden states $(h_1^{(2)}, \dots, h_T^{(2)})$ from the top BiGRU layer, we obtain a fixed-size molecular representation via attention-based pooling. Specifically, we compute intermediate vectors
\begin{equation}
  u_t = \tanh(W_u h_t^{(2)} + b_u), \quad u_t \in \mathbb{R}^{d_a},
\end{equation}
attention weights
\begin{equation}
  \alpha_t = \frac{\exp(u_t^{\top} v)}{\sum_{k=1}^{T} \exp(u_k^{\top} v)}, \quad \alpha_t \geq 0,\ \sum_{t=1}^{T} \alpha_t = 1,
\end{equation}
and the molecular representation
\begin{equation}
  h_{\mathrm{mol}} = \sum_{t=1}^{T} \alpha_t h_t^{(2)} \in \mathbb{R}^{256}.
\end{equation}
Here, $W_u \in \mathbb{R}^{d_a \times 256}$, $b_u \in \mathbb{R}^{d_a}$, and $v \in \mathbb{R}^{d_a}$ are learnable parameters, and $d_a$ is the attention dimension (chosen in practice to be comparable to $d_h$). Attention pooling allows the model to focus on SMILES positions that are most informative for $\lambda_{\max}$ (e.g., core chromophore and key substituents), and the attention weights $\alpha_t$ provide token-level importance scores that can be inspected for interpretability.

\paragraph{Dense head, dropout, and loss.}
The pooled molecular representation $h_{\mathrm{mol}}$ is passed through a dense hidden layer with ReLU activation:
\begin{equation}
  z = \phi_{\mathrm{ReLU}}(W_z h_{\mathrm{mol}} + b_z), \quad W_z \in \mathbb{R}^{128 \times 256},\ b_z \in \mathbb{R}^{128},
\end{equation}
followed by a dropout layer with dropout rate $0.2$ applied to $z$ during training. The final output layer is a dense linear layer
\begin{equation}
  \hat{y} = w^{\top} z + b, \quad w \in \mathbb{R}^{128},\ b \in \mathbb{R},
\end{equation}
which produces the predicted $\lambda_{\max}$ (in nm).

Given experimental target values $y$ (measured $\lambda_{\max}$), the model is trained to minimize the mean absolute error (MAE),
\begin{equation}
  \mathcal{L}_{\mathrm{MAE}} = \frac{1}{N} \sum_{i=1}^{N} \left| y_i - \hat{y}_i \right|,
\end{equation}
where $N$ is the number of training molecules.

\subsection{Model overview schematic}

\begin{figure}[H]
  \centering
  \tikzset{
    block/.style={draw, rounded corners, align=center,
                  minimum height=1.0cm, minimum width=4.5cm, fill=gray!5},
    smallblock/.style={draw, rounded corners, align=center,
                       minimum height=0.9cm, minimum width=4.0cm, fill=gray!5},
    arrow/.style={-{Stealth[length=3mm]}, thick},
  }
  \begin{tikzpicture}[node distance=1.4cm]
    \node[smallblock] (smiles) {SMILES string};
    \node[smallblock, below=of smiles] (tokens) {Character tokenization};
    \node[block, below=of tokens] (embed) {Embedding\\$d_{\mathrm{emb}} = 50$};
    \node[block, below=of embed] (bigru1) {BiGRU layer 1\\$128$ units / direction};
    \node[block, below=of bigru1] (bigru2) {BiGRU layer 2\\$128$ units / direction};
    \node[block, below=of bigru2] (pool) {Global Average Pooling\\$h_{\mathrm{mol}} \in \mathbb{R}^{256}$};
    \node[block, below=of pool] (dense) {Dense (ReLU)\\$128$ units, Dropout $p = 0.2$};
    \node[smallblock, below=of dense] (out) {Linear output: $\hat{\lambda}_{\max}$ (nm)};

    \draw[arrow] (smiles) -- (tokens);
    \draw[arrow] (tokens) -- (embed);
    \draw[arrow] (embed) -- (bigru1);
    \draw[arrow] (bigru1) -- (bigru2);
    \draw[arrow] (bigru2) -- (pool);
    \draw[arrow] (pool) -- (dense);
    \draw[arrow] (dense) -- (out);
  \end{tikzpicture}
  \caption{Vertical schematic of the proposed sequence-based deep learning architecture for predicting the UV absorption maximum, $\lambda_{\max}$, of small-molecule UV absorbers from SMILES.}
  \label{fig:uv_model_arch_supp}
\end{figure}

\end{document}
